% GENERATED FILE. Edit canonical lessons, then run npm run generate:books.
% canonical-source-hash: fnv1a64:2ac7dc1c08914968
% canonical-lessons: FA-C385-jamee, FA-C385-ejtemai, FA-C385-nasl, FA-C385-salmand, FA-C385-mazhab

\chapter{Society and generations}
\label{ch:fa-society-and-generations}
\hlchaptermodality{385}

\begin{chapteropening}
\textbf{By the end of this chapter:} \emph{I can talk about society and social life, the generations, elderly people and religion.}

Complete the last lesson of chapter 385: I can talk about society and social life, the generations, elderly people and religion.
\end{chapteropening}

\section[\texorpdfstring{\fa{جامعه} (jâme'e)}{جامعه (jâme'e)}]{\fa{جامعه} (jâme'e) — a society, a community}

\label{lesson:FA-C385-jamee}



\begin{quote}
Before the new one: say the Persian for to succeed, then the Persian for to be defeated.
\end{quote}

\subsection*{You'll want to know: \fa{جامعه}}
\textbf{\fa{جامعه}} (\emph{jâme'e}) — \textquotedblleft{}a society, a community\textquotedblright{}.

\begin{quote}
\fa{جوانها} \fa{آینده} \fa{جامعه} \fa{هستند}.
\end{quote}

\begin{quote}
\emph{javânhâ âyande-ye jâme'e hastand.}
\end{quote}

\begin{quote}
Young people are the future of society.
\end{quote}

Plural \emph{javâme'}, an Arabic broken plural. \emph{Jâme'e-shenâsi} — sociology.

\begin{cousinweb}
\textbf{\fa{جامعه}} comes from Arabic \emph{jāmi'a} (EtymDB), a word usually explained as what gathers people together; in Persian it is a society.
\end{cousinweb}

\begin{grammarlens}[title={What you've built}]
One more word for telling, explaining and describing at B1.
\end{grammarlens}

\subsection*{Your turn}
\begin{itemize}
\raggedright
  \item \emph{Say it:} \emph{jâme'e}
  \item \emph{Say it:} \emph{jâme'e}, once more
  \item \emph{Say it:} a sentence of your own with \emph{jâme'e}
  \item \emph{Recall:} say the Persian for to succeed, then the Persian for to be defeated, then say \emph{jâme'e} again
\end{itemize}

\begin{tcolorbox}[breakable,colback=teal!4,colframe=teal!35!black,title={Before you move on}]
What does \fa{جامعه} mean? (\textquotedblleft{}A society, a community\textquotedblright{}.) Say it once more.
\end{tcolorbox}
\section[\texorpdfstring{\fa{اجتماعی} (ejtemâ'i)}{اجتماعی (ejtemâ'i)}]{\fa{اجتماعی} (ejtemâ'i) — social}

\label{lesson:FA-C385-ejtemai}



\begin{quote}
Before the new one: say the Persian for to be defeated, then the Persian for a society.
\end{quote}

\subsection*{You'll want to know: \fa{اجتماعی}}
\textbf{\fa{اجتماعی}} (\emph{ejtemâ'i}) — \textquotedblleft{}social\textquotedblright{}.

\begin{quote}
\fa{او} \fa{آدم} \fa{خیلی} \fa{اجتماعی} \fa{است}.
\end{quote}

\begin{quote}
\emph{u âdam-e kheyli ejtemâ'i ast.}
\end{quote}

\begin{quote}
He is a very sociable person.
\end{quote}

From \emph{ejtemâ'} \textquotedblleft{}gathering, society\textquotedblright{} + \emph{-i}. \emph{Masâ'el-e ejtemâ'i} — social issues.

\begin{cousinweb}
\textbf{\fa{اجتماعی}} was borrowed from Arabic \emph{ijtimā'ī} \textquotedblleft{}social\textquotedblright{} (EtymDB), the adjective of \emph{ijtimā'} \textquotedblleft{}gathering\textquotedblright{}.
\end{cousinweb}

\begin{grammarlens}[title={What you've built}]
One more word for telling, explaining and describing at B1.
\end{grammarlens}

\subsection*{Your turn}
\begin{itemize}
\raggedright
  \item \emph{Say it:} \emph{ejtemâ'i}
  \item \emph{Say it:} \emph{ejtemâ'i}, once more
  \item \emph{Say it:} a sentence of your own with \emph{ejtemâ'i}
  \item \emph{Recall:} say the Persian for to be defeated, then the Persian for a society, then say \emph{ejtemâ'i} again
\end{itemize}

\begin{tcolorbox}[breakable,colback=teal!4,colframe=teal!35!black,title={Before you move on}]
What does \fa{اجتماعی} mean? (\textquotedblleft{}Social\textquotedblright{}.) Say it once more.
\end{tcolorbox}
\section[\texorpdfstring{\fa{نسل} (nasl)}{نسل (nasl)}]{\fa{نسل} (nasl) — a generation}

\label{lesson:FA-C385-nasl}



\begin{quote}
Before the new one: say the Persian for a society, then the Persian for social.
\end{quote}

\subsection*{You'll want to know: \fa{نسل}}
\textbf{\fa{نسل}} (\emph{nasl}) — \textquotedblleft{}a generation\textquotedblright{}.

\begin{quote}
\fa{نسل} \fa{جوان} \fa{بیشتر} \fa{با} \fa{اینترنت} \fa{آشناست}.
\end{quote}

\begin{quote}
\emph{nasl-e javân bishtar bâ internet âshenâst.}
\end{quote}

\begin{quote}
The young generation knows the internet better.
\end{quote}

Plural \emph{naslhâ}. Also descent: \emph{az nasl-e} … — descended from … .

\begin{cousinweb}
\textbf{\fa{نسل}} was borrowed from Arabic \emph{nasl} \textquotedblleft{}offspring, descendants\textquotedblright{} (EtymDB); in Persian it is a generation.
\end{cousinweb}

\begin{grammarlens}[title={What you've built}]
One more word for telling, explaining and describing at B1.
\end{grammarlens}

\subsection*{Your turn}
\begin{itemize}
\raggedright
  \item \emph{Say it:} \emph{nasl}
  \item \emph{Say it:} \emph{nasl}, once more
  \item \emph{Say it:} a sentence of your own with \emph{nasl}
  \item \emph{Recall:} say the Persian for a society, then the Persian for social, then say \emph{nasl} again
\end{itemize}

\begin{tcolorbox}[breakable,colback=teal!4,colframe=teal!35!black,title={Before you move on}]
What does \fa{نسل} mean? (\textquotedblleft{}A generation\textquotedblright{}.) Say it once more.
\end{tcolorbox}
\section[\texorpdfstring{\fa{سالمند} (sâlmand)}{سالمند (sâlmand)}]{\fa{سالمند} (sâlmand) — elderly; an elderly person}

\label{lesson:FA-C385-salmand}



\begin{quote}
Before the new one: say the Persian for social, then the Persian for a generation.
\end{quote}

\subsection*{You'll want to know: \fa{سالمند}}
\textbf{\fa{سالمند}} (\emph{sâlmand}) — \textquotedblleft{}elderly; an elderly person\textquotedblright{}.

\begin{quote}
\fa{پرستار} \fa{هر} \fa{روز} \fa{به} \fa{سالمندها} \fa{کمک} \fa{میکند}.
\end{quote}

\begin{quote}
\emph{parastâr har ruz be sâlmandhâ komak mikonad.}
\end{quote}

\begin{quote}
The nurse helps the elderly every day.
\end{quote}

\emph{Sâl} \textquotedblleft{}year\textquotedblright{} + \emph{-mand} \textquotedblleft{}having\textquotedblright{}. More respectful than \emph{pir} \textquotedblleft{}old\textquotedblright{}.

\begin{cousinweb}
\textbf{\fa{سالمند}} is native Persian, built from \emph{sâl} \textquotedblleft{}year\textquotedblright{} + the suffix \emph{-mand} \textquotedblleft{}having\textquotedblright{}. An elderly person is rich in years.
\end{cousinweb}

\begin{grammarlens}[title={What you've built}]
One more word for telling, explaining and describing at B1.
\end{grammarlens}

\subsection*{Your turn}
\begin{itemize}
\raggedright
  \item \emph{Say it:} \emph{sâlmand}
  \item \emph{Say it:} \emph{sâlmand}, once more
  \item \emph{Say it:} a sentence of your own with \emph{sâlmand}
  \item \emph{Recall:} say the Persian for social, then the Persian for a generation, then say \emph{sâlmand} again
\end{itemize}

\begin{tcolorbox}[breakable,colback=teal!4,colframe=teal!35!black,title={Before you move on}]
What does \fa{سالمند} mean? (\textquotedblleft{}Elderly; an elderly person\textquotedblright{}.) Say it once more.
\end{tcolorbox}
\section[\texorpdfstring{\fa{مذهب} (mazhab)}{مذهب (mazhab)}]{\fa{مذهب} (mazhab) — a religion, a faith}

\label{lesson:FA-C385-mazhab}



\begin{quote}
Before the new one: say the Persian for a generation, then the Persian for elderly.
\end{quote}

\subsection*{You'll want to know: \fa{مذهب}}
\textbf{\fa{مذهب}} (\emph{mazhab}) — \textquotedblleft{}a religion, a faith\textquotedblright{}.

\begin{quote}
\fa{در} \fa{این} \fa{شهر} \fa{مردم} \fa{مذهبهای} \fa{مختلفی} \fa{دارند}.
\end{quote}

\begin{quote}
\emph{dar in shahr mardom mazhabhâ-ye mokhtalefi dârand.}
\end{quote}

\begin{quote}
In this city people follow different religions.
\end{quote}

Plural \emph{mazâheb}, an Arabic broken plural. Also a branch within a religion.

\begin{cousinweb}
\textbf{\fa{مذهب}} is usually traced to Arabic \emph{madhhab} \textquotedblleft{}a way to go, a school of thought\textquotedblright{}; Persian says the Arabic \emph{dh} as \emph{z}.
\end{cousinweb}

\begin{grammarlens}[title={What you've built}]
One more word for telling, explaining and describing at B1.
\end{grammarlens}

\subsection*{Your turn}
\begin{itemize}
\raggedright
  \item \emph{Say it:} \emph{mazhab}
  \item \emph{Say it:} \emph{mazhab}, once more
  \item \emph{Say it:} a sentence of your own with \emph{mazhab}
  \item \emph{Recall:} say the Persian for a generation, then the Persian for elderly, then say \emph{mazhab} again
\end{itemize}

\begin{tcolorbox}[breakable,colback=teal!4,colframe=teal!35!black,title={Before you move on}]
What does \fa{مذهب} mean? (\textquotedblleft{}A religion, a faith\textquotedblright{}.) Say it once more.
\end{tcolorbox}
